Three-dimensional fluorescence microscopy is a valuable technique for imaging living biological tissue, but the resulting volumes are often compromised by optical blur, anisotropic resolution, and geometric distortion. These limitations are particularly important in zebrafish cardiac imaging, where accurate visualization of the myocardium is essential for both qualitative interpretation and quantitative analysis. Because the structures of interest are small, curved, and embedded in a complex three-dimensional tissue environment, even modest image degradation can reduce the reliability of downstream measurements.

This report presents a deep learning-based restoration pipeline for degraded fluorescence microscopy volumes of zebrafish cardiac tissue. Rather than relying on an explicitly estimated point spread function, the proposed method learns a corrected representation directly from the image data. In this way, the approach combines geometric correction and super-resolution within a single restoration framework, with the goal of recovering images that are structurally more faithful and more suitable for analysis than the raw input volumes.

The motivation for this work comes from the practical limitations of classical deconvolution methods, which often depend on accurate microscope calibration, reliable metadata, and a well-characterized point spread function. In realistic imaging settings, these assumptions are not always satisfied, especially when acquisition conditions vary across samples or experiments. By using a data-driven strategy instead, the proposed pipeline aims to provide a more flexible and robust alternative for improving three-dimensional fluorescence microscopy of zebrafish myocardium.
